%%%%%%%%%%%%%%%%%%%%%%%%%%%%%%%%%%%%%%%%%%%%%%%%%%%%%%%%%%%%
%% LaTeX Document – Rack Configuration Problem (CSPLib prob031)
%%%%%%%%%%%%%%%%%%%%%%%%%%%%%%%%%%%%%%%%%%%%%%%%%%%%%%%%%%%%
\documentclass[11pt,a4paper]{article}

%% ── Packages ──────────────────────────────────────────────
\usepackage[utf8]{inputenc}
\usepackage[T1]{fontenc}
\usepackage{amsmath,amssymb,amsthm}
\usepackage{booktabs}
\usepackage{array}
\usepackage{multirow}
\usepackage{graphicx}
\usepackage{hyperref}
\usepackage[margin=2.5cm]{geometry}
\usepackage{enumitem}
\usepackage{algorithm}
\usepackage{algpseudocode}
\usepackage{xcolor}
\usepackage{tcolorbox}

%% ── Theorem environments ──────────────────────────────────
\newtheorem{definition}{Definition}
\newtheorem{example}{Example}

%% ── Custom commands ───────────────────────────────────────
\newcommand{\setN}{\mathcal{N}}
\newcommand{\setR}{\mathcal{R}}
\newcommand{\setM}{\mathcal{M}}
\newcommand{\setT}{\mathcal{T}}

\title{\textbf{Rack Configuration Problem}\\[4pt]
       \large CSPLib prob031}
\author{}
\date{\today}

\begin{document}
\maketitle

\begin{abstract}
This document provides a comprehensive description of the
\emph{Rack Configuration Problem} (CSPLib prob031~\cite{csplib}).
We present the problem statement, an illustrative example,
the main constraints, and two mathematical programming formulations:
an Integer Linear Programming (ILP) model and a Boolean Satisfiability
(SAT) encoding.  We also describe the benchmark instances available
on CSPLib.
\end{abstract}

\tableofcontents
\newpage

%% ══════════════════════════════════════════════════════════
\section{Problem Description}
\label{sec:description}
%% ══════════════════════════════════════════════════════════

The Rack Configuration Problem was proposed by
Kızıltan and Hnich~\cite{Kiziltan01symmetrybreaking}
and originates from a real-world hardware design scenario.

\subsection{Informal Statement}

A manufacturer must plug a collection of \emph{electronic cards} into
\emph{racks}.  Each rack has a limited number of \emph{connectors} and
a limited \emph{power supply}; each card consumes exactly one connector
and a certain amount of power.
Before any card can be inserted, the rack must be assigned a
\emph{rack model} that determines its connector count, power budget,
and price.
The goal is to assign cards to racks and choose rack models so that:
\begin{enumerate}[label=(\roman*)]
  \item every card is plugged into exactly one rack,
  \item no rack exceeds its connector or power capacity, and
  \item the total price of the chosen rack models is minimised.
\end{enumerate}

\subsection{Formal Notation}

\begin{definition}[Rack Configuration Problem]
An instance of the problem is defined by:
\begin{itemize}
  \item A set of \textbf{rack models}
        $\setM = \{1, 2, \dots, M\}$.
        Each model $m \in \setM$ has:
        \begin{itemize}
          \item maximum power supply $P_m$,
          \item number of connectors $S_m$,
          \item price $C_m$.
        \end{itemize}
  \item A set of \textbf{card types}
        $\setT = \{1, 2, \dots, T\}$.
        Each type $t \in \setT$ has:
        \begin{itemize}
          \item power requirement $p_t$,
          \item demand $d_t$ (the number of cards of type $t$
                that must be plugged).
        \end{itemize}
  \item A set of \textbf{available racks}
        $\setR = \{1, 2, \dots, R\}$.
\end{itemize}
The total number of cards is $N = \sum_{t \in \setT} d_t$.
We seek an assignment of cards to racks and a model selection for each
rack that satisfies all constraints and minimises total cost.
\end{definition}

%% ══════════════════════════════════════════════════════════
\section{Illustrative Example}
\label{sec:example}
%% ══════════════════════════════════════════════════════════

\begin{example}
Consider \textbf{Instance~1} from CSPLib~\cite{csplib}:
\begin{itemize}
  \item $M = 2$ rack models,\quad $T = 4$ card types,\quad $R = 5$ available racks.
\end{itemize}

\medskip
\begin{minipage}[t]{0.45\textwidth}
\centering
\textbf{Rack Models}\\[4pt]
\begin{tabular}{cccc}
\toprule
Model $m$ & $P_m$ & $S_m$ & $C_m$ \\
\midrule
1 & 150 & 8  & 150 \\
2 & 200 & 16 & 200 \\
\bottomrule
\end{tabular}
\end{minipage}
\hfill
\begin{minipage}[t]{0.45\textwidth}
\centering
\textbf{Card Types}\\[4pt]
\begin{tabular}{ccc}
\toprule
Type $t$ & $p_t$ & $d_t$ \\
\midrule
1 & 20 & 10 \\
2 & 40 &  4 \\
3 & 50 &  2 \\
4 & 75 &  1 \\
\bottomrule
\end{tabular}
\end{minipage}

\medskip\noindent
The total number of cards is $N = 10 + 4 + 2 + 1 = 17$,
requiring at least $\lceil 17/16 \rceil = 2$ racks of model~2
(by connectors alone).
\end{example}

\paragraph{Optimal solution.}
Table~\ref{tab:solution1} shows the optimal assignment for this
instance, as reported in~\cite{Kiziltan01symmetrybreaking}.
Only 3 of the 5 available racks are used.

\begin{table}[ht]
\centering
\caption{Optimal solution for Instance~1
  (cost = 550)~\cite{Kiziltan01symmetrybreaking}.}
\label{tab:solution1}
\begin{tabular}{cc*{4}{c}c}
\toprule
Rack & Model & Type~1 & Type~2 & Type~3 & Type~4 & Cost \\
\midrule
1 & 2 & 6 & 2 & 0 & 0 & 200 \\
2 & 2 & 1 & 2 & 2 & 0 & 200 \\
3 & 1 & 3 & 0 & 0 & 1 & 150 \\
4 & -- & 0 & 0 & 0 & 0 & 0 \\
5 & -- & 0 & 0 & 0 & 0 & 0 \\
\midrule
\multicolumn{2}{c}{\textbf{Total}} & 10 & 4 & 2 & 1 & \textbf{550} \\
\bottomrule
\end{tabular}
\end{table}

\paragraph{Constraint verification.}
We check each active rack against the three constraint families:
\begin{itemize}
  \item \textbf{Rack~1 (Model~2, $P_2\!=\!200$, $S_2\!=\!16$):}
        Connectors: $6+2 = 8 \le 16$ \checkmark.\quad
        Power: $6 \times 20 + 2 \times 40 = 120 + 80 = 200 \le 200$ \checkmark.
  \item \textbf{Rack~2 (Model~2):}
        Connectors: $1+2+2 = 5 \le 16$ \checkmark.\quad
        Power: $1 \times 20 + 2 \times 40 + 2 \times 50
              = 20 + 80 + 100 = 200 \le 200$ \checkmark.
  \item \textbf{Rack~3 (Model~1, $P_1\!=\!150$, $S_1\!=\!8$):}
        Connectors: $3+1 = 4 \le 8$ \checkmark.\quad
        Power: $3 \times 20 + 1 \times 75 = 60 + 75 = 135 \le 150$ \checkmark.
  \item \textbf{Demand:}
        Type~1: $6+1+3 = 10 = d_1$ \checkmark,\quad
        Type~2: $2+2+0 = 4 = d_2$ \checkmark,\\
        Type~3: $0+2+0 = 2 = d_3$ \checkmark,\quad
        Type~4: $0+0+1 = 1 = d_4$ \checkmark.
\end{itemize}
All constraints are satisfied.
The total cost is $200 + 200 + 150 = 550$, which is optimal.

\paragraph{Remark.}
Note that both active racks of Model~2 use their power budget
\emph{exactly} ($200/200$), showing that the
\emph{power constraint} is the binding constraint.
The connector capacity ($S_2 = 16$) is far from being exhausted.

%% ══════════════════════════════════════════════════════════
\section{Main Constraints}
\label{sec:constraints}
%% ══════════════════════════════════════════════════════════

The Rack Configuration Problem involves three families of constraints.

\begin{tcolorbox}[colback=blue!5!white, colframe=blue!50!black,
  title={\textbf{Constraint Summary}}]
\begin{enumerate}
  \item \textbf{Demand Constraint:}
        Every card of every type must be plugged into exactly one rack.
        \[
          \forall\, t \in \setT:\quad
          \sum_{i \in \setR} x_{it} = d_t
        \]

  \item \textbf{Connector-Capacity Constraint:}
        The total number of cards in a rack must not exceed the number
        of connectors of its assigned model.
        \[
          \forall\, i \in \setR:\quad
          \sum_{t \in \setT} x_{it}
          \;\le\;
          \sum_{m \in \setM} S_m \cdot y_{im}
        \]

  \item \textbf{Power-Capacity Constraint:}
        The total power consumption of cards in a rack must not exceed
        the power capacity of its assigned model.
        \[
          \forall\, i \in \setR:\quad
          \sum_{t \in \setT} p_t \cdot x_{it}
          \;\le\;
          \sum_{m \in \setM} P_m \cdot y_{im}
        \]
\end{enumerate}
\end{tcolorbox}

\noindent
Additionally, each rack must be assigned \emph{at most one} model
(or remain unused):
\[
  \forall\, i \in \setR:\quad
  \sum_{m \in \setM} y_{im} \;\le\; 1
\]

%% ══════════════════════════════════════════════════════════
\section{ILP Model}
\label{sec:ilp}
%% ══════════════════════════════════════════════════════════

We present a standard Integer Linear Programming formulation.

\subsection{Decision Variables}

\begin{itemize}
  \item $y_{im} \in \{0,1\}$, for $i \in \setR,\; m \in \setM$:\\
        $y_{im} = 1$ iff rack $i$ is assigned model $m$.
  \item $x_{it} \in \mathbb{Z}_{\ge 0}$, for $i \in \setR,\; t \in \setT$:\\
        number of cards of type $t$ placed in rack $i$.
\end{itemize}

\subsection{Objective Function}

Minimise the total cost of all rack models used:
\begin{equation}\label{eq:ilp-obj}
  \min \quad z = \sum_{i \in \setR} \sum_{m \in \setM} C_m \cdot y_{im}
\end{equation}

\subsection{Constraints}

\paragraph{(C1) Model assignment.}
Each rack uses at most one model:
\begin{equation}\label{eq:ilp-model}
  \sum_{m \in \setM} y_{im} \le 1
  \qquad \forall\, i \in \setR
\end{equation}

\paragraph{(C2) Demand.}
All cards of each type are placed:
\begin{equation}\label{eq:ilp-demand}
  \sum_{i \in \setR} x_{it} = d_t
  \qquad \forall\, t \in \setT
\end{equation}

\paragraph{(C3) Connector capacity.}
\begin{equation}\label{eq:ilp-conn}
  \sum_{t \in \setT} x_{it}
  \;\le\;
  \sum_{m \in \setM} S_m \cdot y_{im}
  \qquad \forall\, i \in \setR
\end{equation}

\paragraph{(C4) Power capacity.}
\begin{equation}\label{eq:ilp-power}
  \sum_{t \in \setT} p_t \cdot x_{it}
  \;\le\;
  \sum_{m \in \setM} P_m \cdot y_{im}
  \qquad \forall\, i \in \setR
\end{equation}

\paragraph{(C5) Linking.}
Cards can only be placed in a rack that has a model assigned.
Let $U_t = d_t$ be an upper bound on the number of cards of type~$t$
in any single rack:
\begin{equation}\label{eq:ilp-link}
  x_{it} \le d_t \cdot \sum_{m \in \setM} y_{im}
  \qquad \forall\, i \in \setR,\; t \in \setT
\end{equation}

\paragraph{(C6) Symmetry breaking (optional).}
Racks assigned the same model are interchangeable, leading
to many symmetric solutions.
Following~\cite{Kiziltan01symmetrybreaking}, we add two constraints.

\medskip\noindent
\textbf{(C6a) Model ordering.}
Impose a non-increasing order on the model indices assigned
to consecutive racks:
\begin{equation}\label{eq:ilp-sym-a}
  \sum_{m \in \setM} m \cdot y_{im}
  \;\ge\;
  \sum_{m \in \setM} m \cdot y_{(i+1)m}
  \qquad \forall\, i \in \{1, \dots, R-1\}
\end{equation}
This ensures that racks using ``larger'' models appear first,
eliminating permutations of entire racks.

\medskip\noindent
\textbf{(C6b) Card-count tie-breaking (Model A).}
Whenever two adjacent racks use the \emph{same} model,
we break the remaining symmetry by requiring that the
earlier rack holds at least as many cards in total:
\begin{equation}\label{eq:ilp-sym-b}
  \sum_{m \in \setM} y_{im} \cdot y_{(i+1)m}
  = 1
  \;\Longrightarrow\;
  \sum_{t \in \setT} x_{it}
  \;\ge\;
  \sum_{t \in \setT} x_{(i+1)t}
  \qquad \forall\, i \in \{1, \dots, R-1\}
\end{equation}
In ILP, this conditional constraint can be linearised with a big-$M$:
\[
  \sum_{t} x_{it} - \sum_{t} x_{(i+1)t}
  \;\ge\;
  -\,N \cdot \Bigl(1 - \sum_{m} y_{im} \cdot y_{(i+1)m}\Bigr)
\]
where $N = \sum_t d_t$ and the product $y_{im} \cdot y_{(i+1)m}$
can be linearised via a standard McCormick relaxation.

\paragraph{Variant (Model B).}
An alternative from~\cite{Kiziltan01symmetrybreaking}
(attributed to Van Hentenryck) replaces (C6b) with a simpler
criterion: order only on a single card type
(e.g., type~$1$):
\begin{equation}\label{eq:ilp-sym-b-alt}
  \sum_{m} y_{im} \cdot y_{(i+1)m} = 1
  \;\Longrightarrow\;
  x_{i1} \ge x_{(i+1)1}
  \qquad \forall\, i
\end{equation}
Model~B is cheaper to propagate but breaks fewer symmetries
than Model~A.

\medskip\noindent
\textbf{Variable domains:}
\[
  y_{im} \in \{0,1\}, \qquad
  x_{it} \in \{0, 1, \dots, \min(d_t,\, \max_{m} S_m)\}
\]

\subsection{Model Size}

The ILP has:
\begin{itemize}
  \item $R \cdot M$ binary variables ($y_{im}$),
  \item $R \cdot T$ integer variables ($x_{it}$),
  \item $R + T + 2R + R \cdot T$ constraints
        (before symmetry breaking).
\end{itemize}

%% ══════════════════════════════════════════════════════════
\section{SAT Model}
\label{sec:sat}
%% ══════════════════════════════════════════════════════════

We now describe a propositional (SAT) encoding.
Since the problem involves integer quantities,
we use a \emph{pseudo-Boolean} (PB) or \emph{order-encoding}
approach.

\subsection{Boolean Variables}

\paragraph{Model selection.}
\begin{equation}
  \mathbf{y}_{im} \in \{\textsc{true}, \textsc{false}\}
  \qquad i \in \setR,\; m \in \setM
\end{equation}
$\mathbf{y}_{im}$ is true iff rack $i$ uses model $m$.

\paragraph{Card assignment (order encoding).}
For each rack $i$ and card type $t$, we introduce
\emph{order variables}:
\begin{equation}
  o_{it}^{(k)} \in \{\textsc{true}, \textsc{false}\}
  \qquad k = 1, \dots, \overline{x}_{it}
\end{equation}
where $\overline{x}_{it} = \min(d_t, \max_m S_m)$ is an upper bound.
The meaning is: $o_{it}^{(k)} = \textsc{true}$ iff
$x_{it} \ge k$ (at least $k$ cards of type $t$ are in rack $i$).

The \emph{channelling axioms} enforce the order:
\begin{equation}\label{eq:sat-order}
  o_{it}^{(k)} \;\Rightarrow\; o_{it}^{(k-1)}
  \qquad \forall\, k = 2, \dots, \overline{x}_{it}
\end{equation}
This is encoded as the clause
$(\neg\, o_{it}^{(k)} \lor o_{it}^{(k-1)})$.

\subsection{Constraint Encoding}

\paragraph{(C1) At-most-one model per rack.}
For each rack $i$, we add an at-most-one (AMO) constraint on
$\{\mathbf{y}_{i1}, \dots, \mathbf{y}_{iM}\}$.
Using the pairwise encoding:
\begin{equation}
  \neg\, \mathbf{y}_{im} \lor \neg\, \mathbf{y}_{im'}
  \qquad \forall\, m < m',\; m, m' \in \setM
  \quad (\text{pairwise AMO})
\end{equation}
Alternative encodings (ladder, product, commander) may be
preferred when $M$ is large.

\paragraph{(C2) Demand constraints.}
For each card type $t$, the sum of cards across all racks must
equal~$d_t$. In the order encoding, $x_{it}$ is represented by
$\sum_k o_{it}^{(k)}$.  We need:
\begin{equation}
  \sum_{i \in \setR} \sum_{k=1}^{\overline{x}_{it}} o_{it}^{(k)}
  \;=\; d_t
  \qquad \forall\, t \in \setT
\end{equation}
This cardinality constraint can be encoded using a
\emph{totalizer}~network or \emph{sequential counter}.

\paragraph{(C3) Connector capacity.}
For each rack $i$:
\begin{equation}
  \sum_{t \in \setT} \sum_{k=1}^{\overline{x}_{it}} o_{it}^{(k)}
  \;\le\;
  \sum_{m \in \setM} S_m \cdot \mathbf{y}_{im}
\end{equation}
Since the right-hand side depends on the model selection,
this becomes a \emph{conditional} cardinality constraint.
For each possible model~$m$, we add:
\begin{equation}\label{eq:sat-conn-m}
  \mathbf{y}_{im}
  \;\Rightarrow\;
  \left(
    \sum_{t \in \setT} \sum_{k=1}^{\overline{x}_{it}} o_{it}^{(k)}
    \;\le\; S_m
  \right)
\end{equation}
The implication is encoded by adding indicator clauses:
if $\mathbf{y}_{im}$ is true, then the cardinality bound
$\le S_m$ must hold.

\paragraph{(C4) Power capacity.}
Similarly, for each rack $i$ and model $m$:
\begin{equation}\label{eq:sat-power-m}
  \mathbf{y}_{im}
  \;\Rightarrow\;
  \left(
    \sum_{t \in \setT} p_t \cdot
    \left(\sum_{k=1}^{\overline{x}_{it}} o_{it}^{(k)}\right)
    \;\le\; P_m
  \right)
\end{equation}
This is a \emph{pseudo-Boolean} (PB) constraint with
coefficients~$p_t$, activated conditionally on $\mathbf{y}_{im}$.

\paragraph{(C5) Linking.}
If no model is assigned to rack~$i$, then no cards should be placed:
\begin{equation}
  \neg\, o_{it}^{(1)}
  \;\lor\;
  \bigvee_{m \in \setM} \mathbf{y}_{im}
  \qquad \forall\, i \in \setR,\; t \in \setT
\end{equation}

\paragraph{(C6) Symmetry breaking.}
As in the ILP model (Section~\ref{sec:ilp}), we add
symmetry-breaking constraints following~\cite{Kiziltan01symmetrybreaking}.

\medskip\noindent
\textbf{(C6a) Model ordering.}
Enforce a non-increasing order on rack model indices.
Let $\text{idx}_i = \sum_{m} m \cdot \mathbf{y}_{im}$
denote the model index of rack~$i$.
We require $\text{idx}_i \ge \text{idx}_{i+1}$.
Using the order encoding, for each pair $(i, i{+}1)$ and
for each model index value $v$, we add:
\begin{equation}
  \mathbf{y}_{(i+1),m}
  \;\Rightarrow\;
  \bigvee_{m' \ge m} \mathbf{y}_{i,m'}
  \qquad \forall\, m \in \setM,\;
  \forall\, i \in \{1,\dots,R{-}1\}
\end{equation}
As a clause: $(\neg\, \mathbf{y}_{(i+1),m} \lor
  \mathbf{y}_{i,m} \lor \mathbf{y}_{i,m+1} \lor \cdots
  \lor \mathbf{y}_{i,M})$.

\medskip\noindent
\textbf{(C6b) Card-count tie-breaking.}
When two adjacent racks $i$ and $i{+}1$ share the same model~$m$,
rack~$i$ must hold at least as many total cards as rack~$i{+}1$.
For each model $m$:
\begin{equation}
  (\mathbf{y}_{im} \land \mathbf{y}_{(i+1)m})
  \;\Rightarrow\;
  \left(
    \sum_{t,k} o_{it}^{(k)}
    \;\ge\;
    \sum_{t,k} o_{(i+1)t}^{(k)}
  \right)
\end{equation}
This conditional cardinality constraint can be encoded by
introducing an activation literal for the conjunction
$\mathbf{y}_{im} \land \mathbf{y}_{(i+1)m}$ and
using a sequential counter or totalizer network.

\subsection{Optimisation Approaches}
\label{sec:sat-opt}

Since a pure SAT solver handles \emph{decision} problems
($\exists$ an assignment satisfying all clauses?),
turning the Rack Configuration Problem into an optimisation problem
requires an outer search strategy.
Let $\Phi$ denote the conjunction of all hard clauses
(C1)--(C5) from the previous subsections, and define the
\emph{cost expression}
\begin{equation}\label{eq:cost-expr}
  \mathrm{cost}(\mathbf{y})
  \;=\;
  \sum_{i \in \setR}\;\sum_{m \in \setM} C_m \cdot \mathbf{y}_{im}.
\end{equation}
We seek a satisfying assignment of $\Phi$ that minimises
$\mathrm{cost}(\mathbf{y})$.
Below we describe five complementary approaches.

% ── 5.3.1  Linear Search ──────────────────────────────────
\subsubsection{Linear Search (Top-Down)}
\label{sec:opt-linear}

The simplest strategy starts from a known feasible solution
(or an upper bound $B_0$) and iteratively tightens the bound.

\begin{enumerate}
  \item Solve $\Phi$ (without any cost bound) to obtain a first
        feasible assignment with cost $z^{*}$.
        Set $B \leftarrow z^{*} - 1$.
  \item Add the PB constraint
        $\mathrm{cost}(\mathbf{y}) \le B$ to $\Phi$.
  \item Solve the augmented formula:
    \begin{itemize}
      \item If \textsc{sat}: record the new solution $z^{*}$,
            set $B \leftarrow z^{*} - 1$, and repeat from step~2.
      \item If \textsc{unsat}: the last recorded $z^{*}$ is optimal.
    \end{itemize}
\end{enumerate}

\paragraph{Complexity.}
In the worst case, the number of SAT calls equals the number of
distinct cost values between the initial upper bound and the optimum.
For the Rack Configuration Problem, costs are multiples of
$\gcd(C_1,\dots,C_M)$, so the number of iterations is at most
\[
  \frac{B_0 - z^{*}_{\mathrm{opt}}}{\gcd(C_1,\dots,C_M)} + 1.
\]

\paragraph{Advantage.}
Very easy to implement; each SAT call is independent.

\paragraph{Disadvantage.}
Can be very slow when the gap between the initial and optimal cost
is large, since every intermediate cost level is tested.

% ── 5.3.2  Binary Search ──────────────────────────────────
\subsubsection{Binary Search}
\label{sec:opt-binary}

Binary search reduces the number of SAT calls to
$O(\log(\text{cost range}))$.

\begin{enumerate}
  \item Compute bounds:
        $L \leftarrow 0$ (or a tighter LP lower bound),
        $U \leftarrow \sum_{i \in \setR} \max_{m} C_m$.
  \item While $L < U$:
    \begin{enumerate}
      \item Set $B \leftarrow \lfloor (L + U) / 2 \rfloor$.
      \item Solve $\Phi \;\land\; (\mathrm{cost}(\mathbf{y}) \le B)$.
      \item If \textsc{sat}: set $U \leftarrow z^{*}$
            (actual cost of the found solution).
      \item If \textsc{unsat}: set $L \leftarrow B + 1$.
    \end{enumerate}
  \item Return the best solution found with cost $U$.
\end{enumerate}

\paragraph{Advantage.}
Logarithmic number of SAT calls in the cost range.

\paragraph{Disadvantage.}
Each call is \emph{independent}; learned clauses from one call
cannot help the next (unless using incremental SAT, see below).
Moreover, intermediate SAT calls may be ``wasted'' on
cost bounds far from the optimum.

% ── 5.3.3  Incremental SAT ───────────────────────────────
\subsubsection{Incremental SAT}
\label{sec:opt-incremental}

Modern SAT solvers (MiniSat, CaDiCaL, Glucose, etc.) support
an \emph{incremental} interface that allows:
\begin{itemize}
  \item \textbf{Adding clauses} between consecutive \texttt{solve()}
        calls without restarting the solver.
  \item \textbf{Assumption literals}: solving under a set of
        \emph{temporary} assumptions that can be retracted.
  \item \textbf{Clause reuse}: all conflict clauses learned in
        previous calls remain in the clause database, pruning the
        search space of subsequent calls.
\end{itemize}

\paragraph{Assumption-based cost bounding.}
To encode the cost bound $\mathrm{cost}(\mathbf{y}) \le B$
in a retractable way, we introduce an \emph{activation literal}
$a_B$ and add the implication:
\begin{equation}\label{eq:activation}
  a_B \;\Rightarrow\; (\mathrm{cost}(\mathbf{y}) \le B)
\end{equation}
This is encoded as a set of clauses where every clause
contains $\neg\, a_B$ as an additional literal.
At solve time, we pass $a_B$ as an \emph{assumption}:
\[
  \texttt{solver.solve(assumptions=[}a_B\texttt{])}
\]
If the result is \textsc{unsat}, the solver can also return
the \emph{subset of assumptions} involved in the conflict
(the \emph{failed assumptions}), which can guide the search.

\paragraph{Incremental linear search.}
Combining incrementality with the top-down linear search gives
the most practical approach:

\begin{algorithm}[ht]
\caption{Incremental SAT optimisation for Rack Configuration}
\label{alg:incr-sat}
\begin{algorithmic}[1]
\State \textbf{Input:} formula $\Phi$ (hard clauses)
\State Initialise SAT solver $\mathcal{S}$; add all clauses of $\Phi$
\State $\mathcal{S}.\texttt{solve()}$
  \Comment{Find first feasible solution}
\If{\textsc{unsat}}
  \State \Return \textsc{infeasible}
\EndIf
\State $z^{*} \leftarrow \mathrm{cost}(\mathcal{S}.\texttt{model()})$
\Comment{Record cost of current solution}
\Loop
  \State Encode $\mathrm{cost}(\mathbf{y}) \le z^{*} - 1$
         with activation literal $a$
  \State Add encoding clauses to $\mathcal{S}$
  \State $\textit{result} \leftarrow
         \mathcal{S}.\texttt{solve(assumptions=[}a\texttt{])}$
  \If{$\textit{result} = \textsc{sat}$}
    \State $z^{*} \leftarrow \mathrm{cost}(\mathcal{S}.\texttt{model()})$
    \Comment{Improved solution}
  \Else
    \State \Return $z^{*}$
    \Comment{Optimal}
  \EndIf
\EndLoop
\end{algorithmic}
\end{algorithm}

\paragraph{Incremental binary search.}
The binary-search scheme can also be made incremental.
At each iteration, a \emph{new} activation literal $a_B$ is created
for the current bound~$B$.
\begin{itemize}
  \item If \textsc{sat} with bound $B$: deactivate $a_B$
        (simply stop assuming it) and create a tighter bound.
  \item If \textsc{unsat}: the conflict clauses learned
        under assumption $a_B$ are retained and may help prove
        \textsc{unsat} for future, tighter bounds.
\end{itemize}

\paragraph{Why incrementality matters.}
\begin{enumerate}
  \item \textbf{Clause reuse.}
        Conflict clauses derived during call~$k$ are valid
        for all subsequent calls (they are consequences of the
        \emph{hard} clauses, not the assumptions).
        This can dramatically reduce solving time, especially
        when consecutive bounds are close.

  \item \textbf{Phase saving.}
        The solver's internal variable-assignment heuristic
        (VSIDS, phase saving) is \emph{warm-started} from the
        previous call, often finding a feasible assignment faster.

  \item \textbf{No re-parsing overhead.}
        The formula is loaded once; only a few clauses are added
        per iteration.
\end{enumerate}

\paragraph{Practical note.}
In PySAT, incremental solving is supported via the
\texttt{Solver} class:
\begin{verbatim}
from pysat.solvers import Glucose4
s = Glucose4()
for cl in hard_clauses:
    s.add_clause(cl)
# First solve
s.solve()
best = compute_cost(s.get_model())
while True:
    bound_clauses, activ = encode_cost_le(best - 1)
    for cl in bound_clauses:
        s.add_clause(cl)
    if s.solve(assumptions=[activ]):
        best = compute_cost(s.get_model())
    else:
        break  # optimal
\end{verbatim}

% ── 5.3.4  Weighted MaxSAT ───────────────────────────────
\subsubsection{Weighted MaxSAT}
\label{sec:opt-maxsat}

Instead of iteratively bounding the cost, a \emph{Weighted Partial
MaxSAT} formulation directly encodes the objective.

\paragraph{Formulation.}
The clause set is partitioned into:
\begin{itemize}
  \item \textbf{Hard clauses} $\Phi_H$: all constraint clauses
        (C1)--(C5);
        these must be satisfied.
  \item \textbf{Soft clauses} $\Phi_S$: for each rack~$i$ and
        model~$m$, add the unit clause
        $(\neg\, \mathbf{y}_{im})$ with \emph{weight} $C_m$.
\end{itemize}
The MaxSAT solver seeks an assignment that satisfies all hard
clauses and minimises the total weight of \emph{falsified} soft
clauses.  A soft clause $(\neg\, \mathbf{y}_{im})$ is falsified
exactly when $\mathbf{y}_{im} = \textsc{true}$, i.e., rack~$i$
uses model~$m$, contributing $C_m$ to the cost.  Thus:
\[
  \text{MaxSAT objective}
  = \!\!\sum_{\substack{(i,m):\\ \mathbf{y}_{im} = \textsc{true}}}
    \!\! C_m
  = \mathrm{cost}(\mathbf{y}).
\]

\paragraph{UNSAT-core guided solving.}
State-of-the-art MaxSAT solvers (e.g., RC2, Open-WBO, CASHWMaxSAT)
use \emph{UNSAT-core guided} algorithms:
\begin{enumerate}
  \item Solve the hard clauses plus all soft clauses as
        \emph{assumptions}.
  \item If \textsc{unsat}, extract the UNSAT core
        (a subset of falsified soft clauses).
  \item Relax the core by introducing auxiliary variables
        (e.g., via the Fu--Malik or OLL algorithm)
        and add new cardinality constraints.
  \item Repeat until \textsc{sat}; the solution is optimal.
\end{enumerate}

\paragraph{Advantage.}
No need to explicitly choose cost bounds.
The solver handles the optimisation internally.

\paragraph{Disadvantage.}
Requires a specialised MaxSAT solver; the number of soft clauses
is $R \cdot M$, which can be large.

% ── 5.3.5  Pseudo-Boolean Optimisation ────────────────────
\subsubsection{Pseudo-Boolean (PB) Optimisation}
\label{sec:opt-pb}

A PB solver natively handles constraints of the form
\[
  \sum_{j} a_j \cdot \ell_j \;\le\; b
  \qquad (a_j, b \in \mathbb{Z},\; \ell_j \text{ literals})
\]
as well as a \emph{linear objective} to minimise.

\paragraph{Formulation.}
The entire Rack Configuration Problem can be expressed in
OPB format:

\begin{itemize}
  \item \textbf{Objective:}
        $\min \sum_{i,m} C_m \cdot \mathbf{y}_{im}$.
  \item \textbf{Hard constraints:}
        Each constraint (C1)--(C5) translates directly into
        PB inequalities without the need for auxiliary variables
        or clause encodings.
\end{itemize}

For example, the power-capacity constraint for rack~$i$ becomes:
\[
  \sum_{t \in \setT} p_t \cdot x_{it}
  \;\le\;
  \sum_{m \in \setM} P_m \cdot \mathbf{y}_{im}
  \quad\Longleftrightarrow\quad
  \sum_{t \in \setT} p_t \cdot x_{it}
  - \sum_{m \in \setM} P_m \cdot \mathbf{y}_{im}
  \;\le\; 0
\]
where integer variable $x_{it}$ is represented by its
order-encoding bits $\sum_k o_{it}^{(k)}$.

\paragraph{Advantage.}
Compact formulation; no auxiliary encoding variables needed for
cardinality or PB constraints.
PB solvers (e.g., RoundingSat, Sat4j) can exploit the linear
structure directly.

\paragraph{Disadvantage.}
PB solvers are less mature and less competitive than SAT solvers
on many benchmarks.

% ── 5.3.6  Comparison ────────────────────────────────────
\subsubsection{Comparison of Approaches}

Table~\ref{tab:opt-compare} summarises the key properties of
each approach.

\begin{table}[ht]
\centering
\caption{Comparison of SAT-based optimisation approaches.}
\label{tab:opt-compare}
\resizebox{\textwidth}{!}{%
\begin{tabular}{@{}lcccc@{}}
\toprule
Approach & \# Solver calls & Clause reuse
         & Solver type & Encoding overhead \\
\midrule
Linear search
  & $O(\Delta_{\text{cost}})$ & No  & SAT & PB bound clauses \\
Binary search
  & $O(\log \Delta_{\text{cost}})$ & No  & SAT & PB bound clauses \\
Incremental linear
  & $O(\Delta_{\text{cost}})$ & \textbf{Yes} & SAT (incr.) & PB bound clauses \\
Incremental binary
  & $O(\log \Delta_{\text{cost}})$ & \textbf{Yes} & SAT (incr.) & PB bound clauses \\
Weighted MaxSAT
  & 1 (internal) & Yes & MaxSAT & $R {\cdot} M$ soft clauses \\
PB optimisation
  & 1 (internal) & Yes & PB & None \\
\bottomrule
\end{tabular}%
}
\end{table}

\noindent
Here $\Delta_{\text{cost}}$ denotes the range of feasible cost values
divided by $\gcd(C_1,\dots,C_M)$.
In practice, \textbf{incremental SAT} with linear search often
performs best for moderately sized instances because it combines
the simplicity of iterative tightening with powerful clause reuse.
For larger instances, \textbf{MaxSAT} solvers with UNSAT-core
guided algorithms tend to be more efficient.

\subsection{SAT Model Size}

\begin{itemize}
  \item Model selection: $R \cdot M$ Boolean variables.
  \item Order encoding: $\sum_{i,t} \overline{x}_{it}$ Boolean variables.
        In the worst case this is $R \cdot T \cdot \max_m S_m$.
  \item Clauses: dominated by cardinality/PB constraint encodings.
\end{itemize}

%% ══════════════════════════════════════════════════════════
\section{Dataset}
\label{sec:dataset}
%% ══════════════════════════════════════════════════════════

The benchmark instances are provided by CSPLib~\cite{csplib}
and originate from early work by
Kızıltan and Hnich~\cite{Kiziltan01symmetrybreaking} and
Sabin et al.~\cite{sabin1999optimization}.

\subsection{Instance Format}

Each instance is described by:
\begin{itemize}
  \item \texttt{nbRackModels}: number of rack models ($M$),
  \item \texttt{nbCardTypes}: number of card types ($T$),
  \item \texttt{nbRacks}: number of available racks ($R$),
  \item \texttt{RackModels}: list of tuples
        $\langle P_m, S_m, C_m \rangle$,
  \item \texttt{CardTypes}: list of tuples
        $\langle p_t, d_t \rangle$.
\end{itemize}

\subsection{Benchmark Instances}

Table~\ref{tab:instances} summarises the four benchmark instances.

\begin{table}[ht]
\centering
\caption{Summary of benchmark instances from CSPLib~\cite{csplib}.}
\label{tab:instances}
\small
\resizebox{\textwidth}{!}{%
\begin{tabular}{cccccll}
\toprule
Inst. & $M$ & $T$ & $R$ & $N$
  & Rack models $\langle P, S, C \rangle$ & Card types $\langle p, d \rangle$ \\
\midrule
1 & 2 & 4 & 5 & 17
  & $\langle 150,8,150 \rangle$, $\langle 200,16,200 \rangle$
  & $\langle 20,10 \rangle$, $\langle 40,4 \rangle$, $\langle 50,2 \rangle$, $\langle 75,1 \rangle$ \\[2pt]
2 & 2 & 4 & 10 & 34
  & $\langle 150,8,150 \rangle$, $\langle 200,16,200 \rangle$
  & $\langle 20,20 \rangle$, $\langle 40,8 \rangle$, $\langle 50,4 \rangle$, $\langle 75,2 \rangle$ \\[2pt]
3 & 2 & 6 & 12 & 45
  & $\langle 150,8,150 \rangle$, $\langle 200,16,200 \rangle$
  & \parbox[t]{5.5cm}{$\langle 10,20 \rangle$, $\langle 20,10 \rangle$, $\langle 40,8 \rangle$,\\ $\langle 50,4 \rangle$, $\langle 75,2 \rangle$, $\langle 100,1 \rangle$} \\[8pt]
4 & 6 & 6 & 9 & 25
  & \parbox[t]{5.5cm}{$\langle 50,2,50 \rangle$, $\langle 100,4,100 \rangle$,\\
     $\langle 150,8,150 \rangle$, $\langle 200,16,200 \rangle$,\\
     $\langle 250,32,250 \rangle$, $\langle 300,64,300 \rangle$}
  & \parbox[t]{5.5cm}{$\langle 20,10 \rangle$, $\langle 40,6 \rangle$, $\langle 50,4 \rangle$,\\
     $\langle 75,2 \rangle$, $\langle 100,2 \rangle$, $\langle 150,1 \rangle$} \\[8pt]
\bottomrule
\end{tabular}%
}
\end{table}

\subsection{Sources and Citations}

\begin{itemize}
  \item The problem specification and instances are hosted on
        \textbf{CSPLib} at
        \url{https://www.csplib.org/Problems/prob031/}~\cite{csplib}.
  \item The problem was introduced in the context of symmetry breaking
        by Kızıltan and Hnich~\cite{Kiziltan01symmetrybreaking},
        who proposed constraint-based models with symmetry-breaking
        techniques.
  \item Sabin et al.~\cite{sabin1999optimization} studied
        optimisation methods for constraint resource problems,
        providing early results on this benchmark.
        (Note: the solution reported in~\cite{sabin1999optimization}
        may contain an error regarding the connector capacity;
        see the CSPLib errata.)
  \item The ILOG Solver manual~\cite{Ilog:Solver02} includes
        a related rack configuration example.
\end{itemize}

%% ══════════════════════════════════════════════════════════
\section{Conclusion}
\label{sec:conclusion}
%% ══════════════════════════════════════════════════════════

The Rack Configuration Problem is a combinatorial optimisation problem
with practical relevance in electronic hardware design.
It can be naturally modelled as an ILP (Section~\ref{sec:ilp})
or encoded as a SAT/MaxSAT problem (Section~\ref{sec:sat}).
The CSPLib benchmark provides four instances of increasing difficulty,
ranging from 17 to 45 cards and up to 6 rack models.

%% ── Bibliography ──────────────────────────────────────────
\bibliographystyle{plain}
\bibliography{references/references}

\end{document}
